\subsection{STANDARD GROUPS†}

If \(\mathrm{S}(\mathbf{X}_{1},\mathbf{X}_{2},\ldots,\mathbf{X}_{n})\) is a formal power series with coefficients in A, then, for all \(x_{1},\ldots,x_{r}\) in m, the series \(\mathrm{S}(x_{1},x_{2},\ldots,x_{r})\) is convergent. More precisely, \(m\times m\times\cdots\times m\) is contained in the domain of absolute convergence of S (Differentiable and Analytic Manifolds, R, 4.1.3).

DEFINITION 1. Let \(r\) be an integer \(\geqslant 0\) . A standard group of dimension \(r\) over \(\mathbf{K}\) is a Lie group \(G\) with the following properties:

\begin{enumerate}
\def\labelenumi{(\roman{enumi})}
\item
  the underlying analytic manifold of G is m × m × · · · × m (r factors);
\item
  there exists a formal power series \(\mathbf{F}\) in \(2r\) variables with coefficients in \(\mathbf{A}^r\) , without constant term, such that \(x.y = \mathbf{F}(x,y)\) for all \(x,y\) in \(\mathbf{G}\) .
\end{enumerate}

Then 0.0 = 0 and hence the identity element of G is the origin of m × m · · · × m.
† The results of nos. 3 and 4 and their proofs remain true when the characteristic of K is > 0.

L(G) will be identified with \(K^r\). By § 5, formula (13), the constants of structure of L(G) with respect to the canonical basis belong to A. We shall need, in the same proof to consider the elements of m × ⋯ × m, now as elements of G, now as elements of L(G).

Example. Let \(G = 1 + \mathbf{M}_{n}(m)\) , which is an open subset of \(\mathbf{M}_{n}(\mathbf{K})\) . If \(x \in G\) , then det \(x \in 1 + m\) and hence \(G \subset \mathbf{GL}(n, K)\) . Clearly \(GG \subset G\) . If \(x = 1 + y\) with \(y \in \mathbf{M}_{n}(m)\) , the calculation of the inverse of a matrix proves first that \(x^{-1} \in \mathbf{M}_{n}(\mathbf{A})\) ; if we write \(x^{-1} = 1 + y'\) , then \(y + y' + yy' = 0\) , hence \(y' \in \mathbf{M}_{n}(m)\) and therefore \(x^{-1} \in G\) . Thus G is an open subgroup of \(\mathbf{GL}(n, K)\) . We identify G with \(m^{n^2}\) by means of the mapping \((\delta_{ij} + y_{ij}) \mapsto (y_{ij})\) . Clearly G is a standard group.

THEOREM 4. Let G be a finite-dimensional Lie group. There exists an open subgroup of G isomorphic to a standard group.

By replacing G by a group isomorphic to an open subgroup of G, the problem is reduced to the case where G is an open subset of \(K^{r}\) , with identity element 0 and where the coordinates of the product x.y and the inverse \(x^{[-1]}\) are given by formulae

\begin{enumerate}
\def\labelenumi{(\arabic{enumi})}
\setcounter{enumi}{2}
\tightlist
\item
\end{enumerate}

\[
(x. y) _ {i} = x _ {i} + y _ {i} + \sum_ {| \alpha | \geqslant 1, | \beta | \geqslant 1} c _ {\alpha \beta i} x ^ {\alpha} y ^ {\beta} (i = 1, 2, \dots , r)\tag{3}
\]

\[
(x ^ {[ - 1 ]}) _ {i} = - x _ {i} + \sum_ {| \alpha | > 1} d _ {\alpha i} x ^ {\alpha} \quad (i = 1, 2, \dots , r)\tag{4}
\]

where the series on the right hand side are convergent for \(x, y\) in \(G\) (§ 5, no. 1). Let \(\lambda \in K^*\) and let the group law be transported from \(G\) to \(G' = \lambda G\) by the homothety of ratio \(\lambda\) . For \(x', y'\) in \(G'\) , the product \(x'.y'\) and the inverse \(x'^{(-1)}\) evaluated in \(G'\) have coordinates

\[
(x ^ {\prime}. y ^ {\prime}) _ {i} = x _ {i} ^ {\prime} + y _ {i} ^ {\prime} + \sum_ {| \alpha | \geqslant 1, | \beta | \geqslant 1} c _ {\alpha \beta i} ^ {\prime} x ^ {\prime \alpha} y ^ {\prime \beta} (i = 1, 2, \dots , r)
\]

\[
(x ^ {\prime [ - 1 ]}) _ {i} = - x _ {i} ^ {\prime} + \sum_ {| \alpha | > 1} d _ {\alpha i} ^ {\prime} x ^ {\prime \alpha} \quad (i = 1, 2, \dots , r)
\]

where

\[
c _ {\alpha \beta i} ^ {\prime} = \lambda^ {- | \alpha | - | \beta | + 1} c _ {\alpha \beta i}, \quad d _ {\alpha i} ^ {\prime} = \lambda^ {- | \alpha | + 1} d _ {\alpha i}.
\]

As the series (3) and (4) are convergent, we see that, for \(|\lambda|\) sufficiently large, for all \(\alpha, \beta, i\) ,

\[
\left| c _ {\alpha \beta i} ^ {\prime} \right| \leqslant 1, \quad \left| d _ {\alpha i} ^ {\prime} \right| \leqslant 1
\]

that is \(c_{\alpha\beta i}^{\prime}\in A\) and \(d_{\alpha i}^{\prime}\in A\) ; and on the other hand \(G^{\prime}\supset m\times m\times\cdots\times m\) . Then \(m\times m\times\cdots\times m\) is an open subgroup of \(G^{\prime}\) and is a standard group.

